problem:  
Let \((X,d,\mathfrak{m})\) and \((Y,d_Y,\mathfrak{m}_Y)\) be \(\mathrm{RCD}(K,N)\) and \(\mathrm{RCD}(K',N')\) spaces, with \(K>0\), \(1<N,N'<\infty\), and suppose \(\pi : X \to Y\) is a **1‑Lipschitz map** admitting a measurable disintegration of \(\mathfrak{m} = \int \mathfrak{m}_y\, d\mathfrak{m}_Y(y)\), where each conditional measure \(\mathfrak{m}_y\) is supported on a compact geodesically convex set \(F_y = \pi^{-1}(y)\).  

Assume further that the pullback of Sobolev functions behaves harmonically, i.e., for every \(f \in W^{1,2}(Y)\) with \(\Delta_Y f = 0\), the composition \(f \circ \pi\) lies in \(W^{1,2}(X)\) and satisfies \(\Delta_X(f \circ \pi) = 0\).  
Call such a map a **Sobolev harmonic morphism** between RCD spaces.  

**Question:**  
If every fiber \((F_y, d|_{F_y}, \mathfrak{m}_y)\) satisfies an \(\mathrm{RCD}(0,N_F)\) condition for some \(N_F < N-N'\), must it follow that \(X\) is (metrically measure‑isomorphic to) a **RCD warped product** \(Y \times_\phi F\) over some representative fiber \(F\) and positive warping function \(\phi : Y \to \mathbb{R}^+\) in the sense of Gigli’s warped product construction?  

Alternatively, can there exist a genuinely non‑smooth “harmonic morphism fibration’’ between RCD spaces for which local metric and measure structures of fibers vary in a non‑warped, singular manner, yet the analytic harmonicity property above still holds?  

Why is it a "good" problem:  
This refined version is logically sound and conceptually precise. It directly tests whether analytic harmonicity within the **Sobolev calculus** of RCD spaces enforces **geometric regularity** of fiber structures—an idea that generalizes the classical Fuglede–Ishihara principle to the non‑smooth synthetic setting. The problem investigates whether harmonic morphisms imply a **warped‑product decomposition theorem** under curvature bounds, bridging key analytical and geometric themes:  

- Harmonic functions and Laplacian calculus on non‑smooth spaces (analysis);  
- Structural decomposition and curvature‑dimension geometry (geometry);  
- Stability and rigidity of disintegrations of measures (probability and metric geometry).  

The statement remains simple—does harmonicity force product structure?—yet it reaches deeply into the foundations of RCD theory, Sobolev analysis, and synthetic curvature rigidity.  
A resolution would provide a new pathway toward understanding when analytic “symmetry’’ implies geometric smoothness, sharpening the boundary between **analytic regularity and metric singularity** in modern geometric analysis.